\documentclass[11pt]{article}
\newcommand{\ListaTitulo}{Gabarito 08: Simulado Prova 1 (Unidade 1)}
% Preâmbulo compartilhado: MATF14, Listas de Exercícios 2026
% Usado por ListaNN.tex e GabaritoNN.tex via % Preâmbulo compartilhado: MATF14, Listas de Exercícios 2026
% Usado por ListaNN.tex e GabaritoNN.tex via % Preâmbulo compartilhado: MATF14, Listas de Exercícios 2026
% Usado por ListaNN.tex e GabaritoNN.tex via \input{preamble}

\usepackage[utf8]{inputenc}
\usepackage[T1]{fontenc}
\usepackage[brazil]{babel}
\usepackage{fullpage}
\usepackage{amsmath,amssymb}
\usepackage{enumitem}
\usepackage{xcolor}
\usepackage{fancyhdr}
\usepackage{titlesec}
\usepackage{listings}
\usepackage{hyperref}
\usepackage{booktabs}
\usepackage{graphicx}
\usepackage{tikz}

\definecolor{matf14azul}{HTML}{0D3B54}
\definecolor{matf14dourado}{HTML}{C9971E}

\titleformat{\section}{\normalfont\Large\bfseries\color{matf14azul}}{\thesection}{1em}{}
\titleformat{\subsection}{\normalfont\large\bfseries\color{matf14azul}}{\thesubsection}{1em}{}

\providecommand{\ListaTitulo}{Lista de Exercícios}

\setlength{\headheight}{14pt}
\pagestyle{fancy}
\fancyhf{}
\lhead{\small MATF14: Estatística Econômica I}
\rhead{\small \ListaTitulo}
\cfoot{\thepage}
\renewcommand{\headrulewidth}{0.4pt}

\lstset{
  basicstyle=\ttfamily\small,
  breaklines=true,
  frame=single,
  columns=fullflexible,
  backgroundcolor=\color{gray!8},
  commentstyle=\color{matf14dourado},
  keywordstyle=\color{matf14azul}\bfseries,
  language=R
}

\newcounter{questaocounter}
\newenvironment{questao}[1]{%
  \stepcounter{questaocounter}%
  \bigskip\noindent\textbf{Questão #1.}\ %
}{\par}

\newcommand{\cabecalho}[2]{%
  \begin{center}
    {\Large\bfseries\color{matf14azul} #1}\\[0.3em]
    {\large #2}
  \end{center}
  \vspace{0.5em}
  \noindent\rule{\linewidth}{0.6pt}
  \vspace{0.5em}
}

\newenvironment{solucao}{%
  \par\smallskip\noindent\textit{\color{matf14dourado!80!black}Solução.}\ %
}{\hfill$\blacksquare$\par}


\usepackage[utf8]{inputenc}
\usepackage[T1]{fontenc}
\usepackage[brazil]{babel}
\usepackage{fullpage}
\usepackage{amsmath,amssymb}
\usepackage{enumitem}
\usepackage{xcolor}
\usepackage{fancyhdr}
\usepackage{titlesec}
\usepackage{listings}
\usepackage{hyperref}
\usepackage{booktabs}
\usepackage{graphicx}
\usepackage{tikz}

\definecolor{matf14azul}{HTML}{0D3B54}
\definecolor{matf14dourado}{HTML}{C9971E}

\titleformat{\section}{\normalfont\Large\bfseries\color{matf14azul}}{\thesection}{1em}{}
\titleformat{\subsection}{\normalfont\large\bfseries\color{matf14azul}}{\thesubsection}{1em}{}

\providecommand{\ListaTitulo}{Lista de Exercícios}

\setlength{\headheight}{14pt}
\pagestyle{fancy}
\fancyhf{}
\lhead{\small MATF14: Estatística Econômica I}
\rhead{\small \ListaTitulo}
\cfoot{\thepage}
\renewcommand{\headrulewidth}{0.4pt}

\lstset{
  basicstyle=\ttfamily\small,
  breaklines=true,
  frame=single,
  columns=fullflexible,
  backgroundcolor=\color{gray!8},
  commentstyle=\color{matf14dourado},
  keywordstyle=\color{matf14azul}\bfseries,
  language=R
}

\newcounter{questaocounter}
\newenvironment{questao}[1]{%
  \stepcounter{questaocounter}%
  \bigskip\noindent\textbf{Questão #1.}\ %
}{\par}

\newcommand{\cabecalho}[2]{%
  \begin{center}
    {\Large\bfseries\color{matf14azul} #1}\\[0.3em]
    {\large #2}
  \end{center}
  \vspace{0.5em}
  \noindent\rule{\linewidth}{0.6pt}
  \vspace{0.5em}
}

\newenvironment{solucao}{%
  \par\smallskip\noindent\textit{\color{matf14dourado!80!black}Solução.}\ %
}{\hfill$\blacksquare$\par}


\usepackage[utf8]{inputenc}
\usepackage[T1]{fontenc}
\usepackage[brazil]{babel}
\usepackage{fullpage}
\usepackage{amsmath,amssymb}
\usepackage{enumitem}
\usepackage{xcolor}
\usepackage{fancyhdr}
\usepackage{titlesec}
\usepackage{listings}
\usepackage{hyperref}
\usepackage{booktabs}
\usepackage{graphicx}
\usepackage{tikz}

\definecolor{matf14azul}{HTML}{0D3B54}
\definecolor{matf14dourado}{HTML}{C9971E}

\titleformat{\section}{\normalfont\Large\bfseries\color{matf14azul}}{\thesection}{1em}{}
\titleformat{\subsection}{\normalfont\large\bfseries\color{matf14azul}}{\thesubsection}{1em}{}

\providecommand{\ListaTitulo}{Lista de Exercícios}

\setlength{\headheight}{14pt}
\pagestyle{fancy}
\fancyhf{}
\lhead{\small MATF14: Estatística Econômica I}
\rhead{\small \ListaTitulo}
\cfoot{\thepage}
\renewcommand{\headrulewidth}{0.4pt}

\lstset{
  basicstyle=\ttfamily\small,
  breaklines=true,
  frame=single,
  columns=fullflexible,
  backgroundcolor=\color{gray!8},
  commentstyle=\color{matf14dourado},
  keywordstyle=\color{matf14azul}\bfseries,
  language=R
}

\newcounter{questaocounter}
\newenvironment{questao}[1]{%
  \stepcounter{questaocounter}%
  \bigskip\noindent\textbf{Questão #1.}\ %
}{\par}

\newcommand{\cabecalho}[2]{%
  \begin{center}
    {\Large\bfseries\color{matf14azul} #1}\\[0.3em]
    {\large #2}
  \end{center}
  \vspace{0.5em}
  \noindent\rule{\linewidth}{0.6pt}
  \vspace{0.5em}
}

\newenvironment{solucao}{%
  \par\smallskip\noindent\textit{\color{matf14dourado!80!black}Solução.}\ %
}{\hfill$\blacksquare$\par}


\begin{document}

\cabecalho{Gabarito 08: Simulado}{Revisão geral, no formato e no nível da Prova 1: Estatística Descritiva (Cap. 1, Aulas 01--08)}

\begin{questao}{1} \textbf{(Setor de atividade e porte)}
\begin{solucao}
\begin{enumerate}[label=(\alph*)]
  \item "Setor de atividade" é qualitativa nominal (sem ordem natural entre Comércio, Indústria,
    Serviços). "Porte" é qualitativa ordinal (Pequena $<$ Média $<$ Grande).
  \item
  \begin{center}
  \begin{tabular}{lcc}
  \toprule
  Setor & $f$ & \% \\
  \midrule
  Comércio  & 28 & 70,0 \\
  Indústria & 8  & 20,0 \\
  Serviços  & 4  & 10,0 \\
  \midrule
  Total     & 40 & 100,0 \\
  \bottomrule
  \end{tabular}
  \end{center}
  \item Um gráfico de barras (ordenado por frequência) é a escolha adequada: 3 categorias
    nominais, comparação de magnitude é o objetivo. Um gráfico de pizza também seria aceitável,
    mas com uma fatia dominante (70\%) e duas pequenas, a comparação entre Indústria e Serviços
    fica mais difícil de julgar visualmente em um círculo do que em barras de mesma base.
\end{enumerate}
\end{solucao}
\end{questao}

\begin{questao}{2} \textbf{(Reclamações mensais)}
\begin{solucao}
Dados: $3,5,4,20,6,4,5,7,3,6,5,4$ ($n=12$). Soma $=72$.
\begin{enumerate}[label=(\alph*)]
  \item $\bar x = 72/12 = 6{,}0$. Ordenados: $3,3,4,4,4,5,5,5,6,6,7,20$, mediana
    $=(5+5)/2=5{,}0$ (6º e 7º valores). Moda: $4$ e $5$ aparecem 3 vezes cada, bimodal.
  \item Desvios ao quadrado (usando $\bar x=6$): $(3-6)^2{=}9$ (×2), $(4-6)^2{=}4$ (×3),
    $(5-6)^2{=}1$ (×3), $(6-6)^2{=}0$ (×2), $(7-6)^2{=}1$ (×1), $(20-6)^2{=}196$ (×1). Soma
    $=18+12+3+0+1+196=230$. $\mathrm{Var} = 230/12\approx 19{,}17$; $dp\approx 4{,}38$.
  \item A loja com 20 reclamações puxa a média para cima (de um valor mais próximo de 4--5, sem
    ela, para 6,0 com ela) e infla fortemente o desvio padrão (a maior parcela isolada da soma de
    quadrados, 196 de 230, vem só dessa loja). A mediana (5,0) e a moda (4 e 5) são bem menos
    afetadas, descrevem melhor "a loja típica"; recomenda-se usar a \textbf{mediana} para
    caracterizar a rede, tratando a loja com 20 reclamações como um caso à parte a ser investigado,
    não como representativo.
\end{enumerate}
\end{solucao}
\end{questao}

\begin{questao}{3} \textbf{(Marketing e novos clientes)}
\begin{solucao}
\begin{enumerate}[label=(\alph*)]
  \item $\bar x = 54/8 = 6{,}75$; $\bar y = 153/8 = 19{,}125$. Calculando
    $\sum(x_i-\bar x)(y_i-\bar y) \approx 84{,}25$, $\sum(x_i-\bar x)^2 \approx 27{,}5$,
    $\sum(y_i-\bar y)^2 \approx 258{,}875$, obtém-se
    $$
    r = \frac{84{,}25}{\sqrt{27{,}5}\sqrt{258{,}875}} \approx \frac{84{,}25}{84{,}35} \approx 0{,}999.
    $$
    Correlação linear positiva praticamente perfeita: meses com mais gasto em marketing tendem,
    de forma muito consistente, a captar mais clientes no mês seguinte.
  \item Meses de maior movimento geral da corretora (sazonalidade, lançamento de um novo produto,
    campanha externa) poderiam levar simultaneamente a mais investimento em marketing \emph{e} a
    mais clientes novos, sem que um cause diretamente o outro, uma variável de confusão comum
    explicando os dois.
\end{enumerate}
\end{solucao}
\end{questao}

\begin{questao}{4} \textbf{(Faturamento de microempresas)}
\begin{solucao}
\begin{enumerate}[label=(\alph*)]
  \item Como média (42) $>$ mediana (31), a distribuição é assimétrica à \textbf{direita}: uma
    cauda de empresas com faturamento bem mais alto puxa a média para cima, enquanto a mediana
    permanece próxima da maioria das empresas.
  \item A \textbf{mediana} (R\$ 31 mil): por não ser distorcida pela cauda de faturamentos altos,
    ela descreve melhor a experiência da microempresa "do meio" da distribuição, mais representativa
    do "típico" para um público leigo do que a média, inflada por poucos casos atípicos.
\end{enumerate}
\end{solucao}
\end{questao}

\end{document}
